\noindent\textbf{(a)} \((1-x^2)0-x+x=0\).

\noindent\textbf{(b)} Puede tomarse \(y_2=\sqrt{1-x^2}\) en \((-1,1)\), y \(y_2=\sqrt{x^2-1}\) en \((1,\infty)\) o \((-\infty,-1)\).

\noindent\textbf{(c)} \(y=C_1x+C_2\sqrt{|1-x^2|}\) en cada uno de los tres intervalos regulares. Para prolongar clásicamente a través de \(x=\pm1\), debe ser \(C_2=0\) en las ramas adyacentes.
